\respcap{cap:avl}
\resp{avl-r1} (a) Uma ABB é balanceada se sua altura é $O(\log n)$ e o mesmo vale para cada subárvore (com $m$ nós, altura $O(\log m)$). (b) AVL: em todo nó, as alturas das subárvores esquerda e direita diferem de no máximo 1. (c) Toda AVL é balanceada, pois $h<1{,}44\log_2(n+2)$ (Teorema~\ref{teo:avlaltura}); nem toda árvore balanceada é AVL. (d) $\bal(v)=h_D(v)-h_E(v)$.

\resp{avl-r2} (a) $\bal(p)=1\to0$: a altura não muda e o processo para. $0\to-1$: a altura aumenta e o processo continua com o pai. $-1\to-2$: rotação, e o processo para. (b) Com sinais iguais ($-2$ e $-1$), rotação \textbf{simples à direita} em $p$. Com sinais opostos ($-2$ e $+1$), rotação \textbf{dupla à direita}: à esquerda em $u$ e depois à direita em $p$. No primeiro caso o excesso está ``em linha'' (esquerda-esquerda) e uma rotação resolve; no segundo, está em zigue-zague (esquerda-direita), e é preciso primeiro alinhar.

\resp{avl-r4} (a) $h$ indica se a altura da subárvore aumentou. (b) Se a altura não mudou, nenhum ancestral muda de balanço; só quando $h$ é verdadeiro o pai precisa se atualizar. (c) A rotação muda a raiz da subárvore; como $pt$ é passado por referência, atribuir a nova raiz a $pt$ religa a subárvore ao pai (ou atualiza a raiz da árvore). (d) Os testes precisam ser exclusivos: com ``se'' em sequência, um nó com balanço 1 vira 0 e, logo em seguida, o teste ``balanço $=0$'' dispara e o deixa com $-1$, com $h$ verdadeiro, o que está errado.

\resp{avl-r5} Ver §\ref{sec:avlexemplos}. (a) Rotação simples à direita em 40 ao inserir 10, e dupla à direita em 40 ao inserir 35. Árvore final: $20^{+1}(10,\,35(30,40))$. (b) Rotação simples à esquerda em 10 (ao inserir 30) e simples à esquerda em \textbf{30} (ao inserir 50; o primeiro desregulado é 30, não 20). Árvore final: $20^{+1}(10,\,40(30,50))$. (c) Inserir 49 desregula 20, cujo filho direito 40 tem balanço $+1$: rotação simples à esquerda em 20. Árvore final: $40(20(10,30),\,50^{-1}(49,\cdot))$.

\resp{avl-t5} Ver o Exemplo~\ref{ex:teste1q5}. Inserir 5 desregula 30 (caso esquerda-esquerda, rotação simples à direita em 30); 15 não causa rotação; 35 desregula 50 (caso esquerda-direita, rotação dupla à direita: à esquerda em 20, à direita em 50, e 30 vira a raiz); 45 não causa rotação. Árvore final: $30(20(10(5,15),25),\allowbreak\,50(40(35,45),70(60,80)))$. (c) Numa ABB, inserções em ordem crescente, decrescente ou em zigue-zague geram altura $\Theta(n)$ e operações $\Theta(n)$; a AVL garante $O(\log n)$ qualquer que seja a ordem de chegada, com $O(1)$ rotações por inserção.

\resp{avl-l3} Use uma árvore de Fibonacci (uma AVL mínima) em que todos os nós internos têm balanço $-1$ e remova a folha mais à direita. Cada ancestral perde altura do lado que já era mais baixo, fica com balanço $-2$ e tem o filho esquerdo com balanço $-1$; a rotação simples diminui a altura da subárvore e propaga o problema para cima. Por exemplo, na árvore de Fibonacci de altura 7, com 33 nós (chaves $1,\dots,33$) e raiz 21, remover a folha 33 (cujos ancestrais são 32, 29 e 21) causa rotações em 32, 29 e 21.

\resp{avl-l4} Se todas as chaves de $T_1$ são menores que as de $T_2$ (junção): retire o mínimo $x$ de $T_2$; supondo $h(T_1)\ge h(T_2)$, desça pela espinha direita de $T_1$ até um nó $v$ de altura $h(T_2)$ ou $h(T_2)+1$; crie um nó $x$ com filhos $T(v)$ e $T_2$, coloque-o no lugar de $v$ e suba rebalanceando como numa inserção. Custo $O(\log n)$. No caso geral, com chaves intercaladas: faça os percursos em ordem das duas árvores ($O(n+m)$), intercale as duas listas ordenadas e construa uma árvore completa tomando sempre a mediana como raiz, em $O(n+m)$ no total. Alternativa: inserir os nós da menor na maior, em $O(m\log(n+m))$.

\resp{avl-l5} Sim: cada conjunto é uma AVL e cada nó guarda um ponteiro para o pai. \Find$(x)$ sobe até a raiz da AVL de $x$, que identifica o conjunto: $O(\log n)$. \Union{} insere os elementos da AVL menor na maior: $O(k\log n)$, sendo $k$ o tamanho da menor; com ``menor na maior'', cada elemento é reinserido no máximo $\log n$ vezes, e o total é $O(n\log^2n)$. Ganha-se ordem dentro de cada conjunto, mas a estrutura é bem pior que o Union-Find, de custo praticamente $O(1)$ amortizado.

\resp{avl-r6} (a) AVL: a ABB pode degenerar para altura $\Theta(n)$ com inserções em ordem, enquanto a AVL garante $O(\log n)$ para busca, inserção e remoção. (b) Heap: o máximo está na raiz ($O(1)$) e inserir ou remover custa $O(\log n)$, num vetor simples. A AVL também faria tudo em $O(\log n)$, mas é mais pesada, e só compensaria se também fosse preciso buscar ou remover tarefas específicas por identificador.

\resp{avl-pseudo} A inserção e as rotações estão na §\ref{sec:avl}. Remoção: remova como numa ABB (com dois filhos, troque pelo maior da subárvore esquerda), suba a partir do nó fisicamente removido usando a tabela ``a altura diminuiu?'' e, ao rotacionar, trate o caso em que o filho do lado mais pesado tem balanço 0 com rotação \emph{simples}; nesse caso a altura não muda e o processo para.

\resp{avl-juiz} 1195: insira as chaves numa ABB e imprima os três percursos. 1200 e 1201: ABB com inserção, busca, percursos e, no segundo, remoção. 1191: o primeiro símbolo da pré-ordem é a raiz; sua posição no percurso em ordem separa a subárvore esquerda (à esquerda dela) da direita; resolva recursivamente as duas partes e imprima a raiz por último. Custa $O(n^2)$ na forma simples, ou $O(n)$ com uma tabela de posições.

\resp{avl-rn} Ver §\ref{sec:rn}: definição (raiz e folhas nulas pretas, nunca dois vermelhos seguidos, mesma altura negra em todos os caminhos); $h\le2\log_2(n+1)$, porque a subárvore de $x$ tem pelo menos $2^{bh(x)}-1$ nós internos e pelo menos metade dos nós de um caminho é preta; inserção com no máximo duas rotações (casos do tio vermelho e do tio preto) e remoção com no máximo três.

\resp{avl-1a7} Há rotação simples à esquerda em 1 (ao inserir 3), em 3 (ao inserir 5), em 2 (ao inserir 6; 2 é a raiz nesse momento) e em 5 (ao inserir 7). A árvore final é perfeita: $4(2(1,3),\,6(5,7))$, com todos os balanços iguais a 0.

\resp{avl-3rot} Inserir 30: rotação simples à esquerda em 10, dando $20(10,30)$. Inserir 25 e 28: o nó 30 fica com $-2$ e o filho 25 com $+1$; rotação dupla à direita em 30, dando $20(10,\,28(25,30))$. Inserir 27 (à direita de 25): o nó 20 fica com $+2$ e o filho 28 com $-1$; rotação dupla à esquerda em 20. Árvore final: $25^{0}(20^{-1}(10,\cdot),\ 28^{0}(27,30))$.

\resp{avl-rem} Remover 1: $4(2^{+1}(\cdot,3),\,6(5,7))$, sem rotação. Remover 3: $4^{+1}(2,\,6(5,7))$, sem rotação. Remover 2: o nó 4 fica com $+2$ e o filho 6 com balanço 0, caso exclusivo da remoção; rotação simples à esquerda em 4, dando $6^{-1}(4^{+1}(\cdot,5),\,7)$, e a altura não diminuiu.

\resp{avl-altmax} $N(1),\dots,N(7)=1,2,4,7,12,20,33$. Com 20 nós cabe uma AVL mínima de altura 6 (uma árvore de Fibonacci), e a altura 7 exigiria 33 nós: a altura máxima é 6. Com 32 nós, ainda $N(7)=33>32$: a altura máxima continua 6.

\resp{avl-inordem} Na rotação simples à direita em $p$, com filho esquerdo $u$: antes, o percurso em ordem é $T_1,u,T_2,p,T_3$; depois, com $u$ na raiz, $T_1$ à esquerda e $p$ à direita com $T_2$ e $T_3$, ele é $T_1,u,T_2,p,T_3$ --- o mesmo. A rotação à esquerda é simétrica, e a dupla é a composição de duas simples. Como uma árvore binária é ABB exatamente quando o percurso em ordem é crescente, a propriedade se mantém.

\resp{avl-catalan} $C_3=5$ e $C_4=14$ (números de Catalan). Escolhendo a raiz $r$, a subárvore esquerda guarda $r-1$ chaves e a direita $n-r$, logo $C_n=\sum_{r=1}^nC_{r-1}C_{n-r}$, com $C_0=1$. Os valores são $1,1,2,5,14,42,132$.

\resp{avl-rem30} O maior elemento da subárvore esquerda é 25, que toma o lugar de 30 e é retirado da sua posição. O nó 20 fica sem filho direito e com a subárvore esquerda $10(5,15)$ de altura 2: $\bal(20)=-2$, com o filho 10 de balanço 0. Rotação simples à direita em 20 (caso exclusivo da remoção): a subárvore passa a ser $10^{+1}(5,\,20^{-1}(15,\cdot))$, sem diminuir de altura, e o processo para. Árvore final: $25(10(5,20(15,\cdot)),\,50(40(35,45),70(60,80)))$.

\resp{avl-seq} Rotações: ao inserir 1, simples à direita em 3; ao inserir 5, simples à esquerda em 3; ao inserir 6, simples à esquerda em 2 (a raiz); ao inserir 7, simples à esquerda em 5; ao inserir 15, dupla à esquerda em 7; ao inserir 14, dupla à esquerda em 6. As inserções de 3, 2, 4 e 16 não causam rotação. Árvore final: $4^{+1}(2(1,3),\ 7(6^{-1}(5,\cdot),\,15(14,16)))$.

\resp{avl-hmin} Uma árvore binária de altura $h$ tem no máximo $1+2+\cdots+2^{h-1}=2^h-1$ nós; logo $n\le2^h-1$, isto é, $h\ge\log_2(n+1)$, e, como $h$ é inteiro, $h\ge\ceil{\log_2(n+1)}$. Com o Teorema~\ref{teo:avlaltura}, a altura de uma AVL fica entre $\log_2(n+1)$ e $1{,}44\log_2(n+2)$: é $\Theta(\log n)$.

\resp{avl-verif} Uma função recursiva que, para cada nó, devolve (é válida, altura, menor chave, maior chave) da subárvore: a subárvore de $v$ é válida se as dos filhos são válidas, se a maior chave da esquerda é $\le v.chave\le$ a menor chave da direita e se as alturas dos filhos diferem de no máximo 1. Cada nó é visitado uma vez: $O(n)$.

\resp{avl-suc} Se $x$ tem filho direito, o sucessor é o mínimo da subárvore direita (desça uma vez à direita e depois sempre à esquerda). Senão, suba pelos pais enquanto $x$ for filho direito; o sucessor é o primeiro ancestral do qual $x$ está na subárvore esquerda (ou não existe, se chegar à raiz por um caminho só de filhos direitos). Em ambos os casos o caminho percorrido tem comprimento no máximo $h$.
